%
% client_calls_replay_on_reconnect_buggy.tex
%
% FIDELITY VARIANT of docs/client_calls.tex.  Pending calls survive a session
% end and are re-sent to the successor: the callback hold's behaviour applied
% to requests.  A call executes on a session it was not sent to.  C4 must fail.
%
% A Z specification of the Hub's typed request to a provider client and its
% reply (DES-107): the registration map, the table of pending calls, and the
% five outcomes a call can have.
%
% Design record: DESIGN.md, DES-107 ("The Hub Relays a Typed Request to a
% Provider Client and Returns Its Reply"), decisions 1 to 4.
% Builds on DES-102 (a reply reaches only the dial that asked: here, the
% listen session that was asked), DES-101 (registration and withdrawal are
% lifecycle steps) and DES-058 (the callback hold, which requests must not
% copy).
%
% The lock order is checked in docs/client_calls_locks.tex.
%
% Type-check:  fuzz -t docs/client_calls_replay_on_reconnect_buggy.tex
% Model-check: scripts/verify_client_calls_model.sh
%
\documentclass[a4paper,10pt,fleqn]{article}
\usepackage[margin=1in]{geometry}
\usepackage{fuzz}

\begin{document}

\title{The lux Client Calls: a Z Specification}
\author{Formal model of the Hub-to-provider request relay}
\date{November 2026}
\maketitle

\section{Introduction}

A caller on a Hub worker thread asks a provider client a question.  The
request goes down the provider's listen leg as one frame; one frame comes back.
Between the two, the Hub holds a table of pending calls.  The question the
model answers is the one DES-102 answered for the Display dial: who may
complete a call, and what happens to a call when the thing it was sent to
goes away.

Every call has exactly one outcome.  The model makes that a checked
property, not a convention.

\section{Basic Types}

\begin{zed}
  [CONN, REQ, KEY]
\end{zed}

\texttt{CONN} is a listen session: one live listen leg of one client, from
its handshake to its end.  A client that reconnects has a new session.
Session identities are never reused, so a reply cannot be credited to a
successor of the session it came from.  \texttt{REQ} is a request id
(\texttt{c<boot>.<n>}), minted by the Hub and never reused.  \texttt{KEY} is
a registration key, the pair (service, host).

\begin{zed}
  OUT ::= replied | refused | timedout | lost
\end{zed}

The four outcomes of a call that was sent.  A call that was never sent (no
provider, provider busy, the caller not identified) has no entry in the
table and is not modelled.

\begin{axdef}
  maxPend : \nat
\where
  maxPend = 2
\end{axdef}

The per-provider limit.  The design says 8; the model scales it to 2 so that
a carrier of three requests reaches it.

\section{State}

\begin{schema}{CC}
  live : \power CONN \\
  opened : \power CONN \\
  reg : KEY \pfun CONN \\
  pend : \power REQ \\
  sentTo : REQ \pfun CONN \\
  inflight : \power (REQ \cross CONN) \\
  served : \power (REQ \cross CONN) \\
  outc : REQ \pfun OUT \\
  accFrom : REQ \pfun CONN \\
  nout : REQ \fun 0 \upto 2
\where
  live \subseteq opened \\
  pend \subseteq \dom sentTo \\
  \dom outc \subseteq \dom sentTo \\
  \dom accFrom \subseteq \dom outc
\end{schema}

\texttt{live} are the sessions now connected; \texttt{opened} every session
ever started.  \texttt{reg} is the registration map: the session that serves
each key.  Being a function it gives at most one session per key.
\texttt{pend} is the table of pending calls; \texttt{sentTo} is the history
of which session each request was sent to, and is never rewritten.
\texttt{inflight} are request frames queued for a session, \texttt{served}
those a provider has taken up and run.  \texttt{outc} is the outcome the
waiter was given, \texttt{accFrom} the session a reply or refusal was
accepted from, and \texttt{nout} the number of outcomes delivered to each
request's waiter.

\begin{schema}{Init}
  CC'
\where
  live' = \emptyset \\
  opened' = \emptyset \\
  reg' = \emptyset \\
  pend' = \emptyset \\
  sentTo' = \emptyset \\
  inflight' = \emptyset \\
  served' = \emptyset \\
  outc' = \emptyset \\
  accFrom' = \emptyset \\
  nout' = REQ \cross \{ 0 \}
\end{schema}

\section{Session Operations}

\begin{schema}{OpenConn}
  \Delta CC \\
  c? : CONN
\where
  c? \notin opened \\
  live' = live \cup \{ c? \} \\
  opened' = opened \cup \{ c? \} \\
  reg' = reg \\
  pend' = pend \\
  sentTo' = sentTo \\
  inflight' = inflight \\
  served' = served \\
  outc' = outc \\
  accFrom' = accFrom \\
  nout' = nout
\end{schema}

A client's listen leg completes its handshake.

\begin{schema}{SessionEnd}
  \Delta CC \\
  c? : CONN \\
  dead! : \power REQ
\where
  c? \in live \\
  dead! = \{ r : \dom sentTo | r \in pend \land sentTo~r = c? \spot r \} \\
  live' = live \setminus \{ c? \} \\
  opened' = opened \\
  reg' = reg \nrres \{ c? \} \\
  pend' = pend \setminus dead! \\
  sentTo' = sentTo \\
  inflight' = inflight \nrres \{ c? \} \\
  served' = served \\
  outc' = outc \oplus (dead! \cross \{ lost \}) \\
  accFrom' = accFrom \\
  nout' = nout \oplus
    \{ r : REQ | r \in dead! \spot r \mapsto nout~r + 1 \}
\end{schema}

The session ends: the client disconnected, its lease lapsed, or luxd is
tearing down.  One step withdraws the session's registrations, discards the
frames queued for it, and resolves every call pending on it as
\texttt{lost}.  This is C4: no pending call outlives its session, and nothing
is sent to a successor.

\begin{schema}{Supersede}
  \Delta CC \\
  c? : CONN \\
  n? : CONN \\
  dead! : \power REQ
\where
  c? \in live \\
  n? \notin opened \\
  dead! = \{ r : \dom sentTo | r \in pend \land sentTo~r = c? \spot r \} \\
  live' = (live \setminus \{ c? \}) \cup \{ n? \} \\
  opened' = opened \cup \{ n? \} \\
  reg' = reg \nrres \{ c? \} \\
  pend' = pend \\
  sentTo' = sentTo \\
  inflight' = (inflight \nrres \{ c? \}) \cup
    \{ r : REQ | r \in dead! \spot (r , n?) \} \\
  served' = served \\
  outc' = outc \\
  accFrom' = accFrom \\
  nout' = nout
\end{schema}

A newer listen leg of the same client replaces the old.  The defect: the
calls pending on the old session survive and are queued again for the new
one, the callback hold's redelivery applied to requests.  The caller was
never told \texttt{lost}, yet the request runs on a session it was not sent
to.  C4 fails, both parts.

\section{Registration Operations}

\begin{schema}{RegisterSvc}
  \Delta CC \\
  c? : CONN \\
  k? : KEY
\where
  c? \in live \\
  k? \notin \dom reg \lor reg~k? = c? \\
  reg' = reg \oplus \{ k? \mapsto c? \} \\
  live' = live \\
  opened' = opened \\
  pend' = pend \\
  sentTo' = sentTo \\
  inflight' = inflight \\
  served' = served \\
  outc' = outc \\
  accFrom' = accFrom \\
  nout' = nout
\end{schema}

A key is free, or held by the connection that asks.  The same connection
registering again, as a restarted helper does, replaces its own entry.  The
guard \texttt{c? $\in$ live} is the listen-leg gate (\texttt{push\_required}).

\begin{schema}{RegisterRefused}
  \Xi CC \\
  c? : CONN \\
  k? : KEY
\where
  c? \in live \\
  k? \in \dom reg \\
  reg~k? \neq c?
\end{schema}

A different connection claiming a held key is refused and nothing changes.

\section{Call Operations}

\begin{schema}{Dispatch}
  \Delta CC \\
  r? : REQ \\
  k? : KEY
\where
  k? \in \dom reg \\
  r? \notin \dom sentTo \\
  \# \{ q : \dom sentTo | q \in pend \land sentTo~q = reg~k? \spot q \}
    < maxPend \\
  pend' = pend \cup \{ r? \} \\
  sentTo' = sentTo \cup \{ r? \mapsto reg~k? \} \\
  inflight' = inflight \cup \{ (r? , reg~k?) \} \\
  live' = live \\
  opened' = opened \\
  reg' = reg \\
  served' = served \\
  outc' = outc \\
  accFrom' = accFrom \\
  nout' = nout
\end{schema}

A caller mints a fresh id, enters the call as pending for the provider's
current session and queues the frame for that session.  The ninth call
(here the third) is refused at once and never enters the table; a refusal
that changes nothing is not an operation of the model.

\begin{schema}{Serve}
  \Delta CC \\
  r? : REQ \\
  c? : CONN
\where
  (r? , c?) \in inflight \\
  c? \in live \\
  inflight' = inflight \setminus \{ (r? , c?) \} \\
  served' = served \cup \{ (r? , c?) \} \\
  live' = live \\
  opened' = opened \\
  reg' = reg \\
  pend' = pend \\
  sentTo' = sentTo \\
  outc' = outc \\
  accFrom' = accFrom \\
  nout' = nout
\end{schema}

The provider takes the frame from its leg and runs the request.  Serving is
not guarded by the table: a call that has already timed out may still be
served, which is why a timeout reports an unknown result.

\begin{schema}{ReplyFrame}
  \Delta CC \\
  r? : REQ \\
  c? : CONN
\where
  c? \in live \\
  r? \in pend \\
  sentTo~r? = c? \\
  pend' = pend \setminus \{ r? \} \\
  outc' = outc \oplus \{ r? \mapsto replied \} \\
  accFrom' = accFrom \oplus \{ r? \mapsto c? \} \\
  nout' = nout \oplus \{ r? \mapsto nout~r? + 1 \} \\
  live' = live \\
  opened' = opened \\
  reg' = reg \\
  sentTo' = sentTo \\
  inflight' = inflight \\
  served' = served
\end{schema}

A reply frame arrives on live session \texttt{c?} naming \texttt{r?}.  It
counts only if the request is pending and was sent to that session.  This
is C2.  The session \texttt{c?} is any live connection, the provider or not:
the model lets the environment send any frame from any live leg.  A frame
that fails the guard is dropped, and a dropped frame changes nothing, so it
is not an operation either.  Dropping covers C3: a late frame for a call
that has already had its outcome finds no entry.

\begin{schema}{RefuseFrame}
  \Delta CC \\
  r? : REQ \\
  c? : CONN
\where
  c? \in live \\
  r? \in pend \\
  sentTo~r? = c? \\
  pend' = pend \setminus \{ r? \} \\
  outc' = outc \oplus \{ r? \mapsto refused \} \\
  accFrom' = accFrom \oplus \{ r? \mapsto c? \} \\
  nout' = nout \oplus \{ r? \mapsto nout~r? + 1 \} \\
  live' = live \\
  opened' = opened \\
  reg' = reg \\
  sentTo' = sentTo \\
  inflight' = inflight \\
  served' = served
\end{schema}

A refusal frame obeys the rule a reply obeys.

\begin{schema}{Timeout}
  \Delta CC \\
  r? : REQ
\where
  r? \in pend \\
  pend' = pend \setminus \{ r? \} \\
  outc' = outc \oplus \{ r? \mapsto timedout \} \\
  nout' = nout \oplus \{ r? \mapsto nout~r? + 1 \} \\
  live' = live \\
  opened' = opened \\
  reg' = reg \\
  sentTo' = sentTo \\
  inflight' = inflight \\
  served' = served \\
  accFrom' = accFrom
\end{schema}

The deadline is abstracted to a guard: the call is pending, and the
transition may fire at any later point.  This is C5.  Without it a provider
that never answers pins the waiter, and one of the provider's slots, for
ever.

\begin{schema}{Quiescent}
  \Xi CC
\where
  live = \emptyset \\
  opened = CONN
\end{schema}

The finite carrier has run out of sessions and all of them have ended.  It
is enabled only then, so that the deadlock check reports any other state in
which nothing can move.

\section{Invariants}
\label{sec:inv}

None is placed in \texttt{CC}: a predicate there would disable the very
transition that breaks it.  Each is checked as an unreachable goal.

\paragraph{C1 --- a call leaves pending once.}
No waiter is given two outcomes; a pending call has been given none; a call
that is no longer pending has been given one:
\[
  \forall r : REQ @ nout~r \leq 1 \land (r \in pend \implies nout~r = 0)
    \land (r \in \dom sentTo \setminus pend \implies nout~r = 1)
\]

\paragraph{C2 --- a reply counts only from the session it was sent to.}
\[
  \forall r : \dom accFrom @ accFrom~r = sentTo~r
\]

\paragraph{C3 --- an outcome removes the entry.}
\[
  \dom outc \cap pend = \emptyset
\]

\paragraph{C4 --- a session's end resolves its calls in the same step.}
Two parts.  C4a, no pending call outlives its session: every pending call
is on a live session.  C4b, nothing is redelivered: a request is served only
by the session it was sent to.
\[
  \forall r : pend @ sentTo~r \in live
\]
\[
  \forall (r , c) \in served @ sentTo~r = c
\]

\paragraph{C5 --- a call that reaches the deadline times out.}
Whenever a call is pending, \texttt{Timeout} is enabled.  It is checked as the
CTL formula $\mathrm{AG}(P \Rightarrow \mathrm{EF}\,\mathrm{enabled}(
\mathit{Timeout}))$, where $P$ says that some call is pending.  $P$ is
reachable, so the formula is not vacuous.

\paragraph{C6 --- no lock is held while waiting.}
Checked in \texttt{client\_calls\_locks.tex}.

\paragraph{R1 --- a registration is held by a live session.}
\[
  \forall k : \dom reg @ reg~k \in live
\]

\paragraph{P1 --- the per-provider limit holds.}
\[
  \forall c : CONN @ \# \{ r : pend | sentTo~r = c \} \leq maxPend
\]

\paragraph{D --- deadlock freedom.}
Some operation is always enabled, until the carrier is spent.

\section{Verification}
\label{sec:verify}

\begin{verbatim}
  scripts/verify_client_calls_model.sh
\end{verbatim}

\section{Fidelity}

Four buggy variants sit beside this file; each must make a goal reachable
that the intact model cannot reach (see
\texttt{scripts/verify\_client\_calls\_model.sh} for the exact commands).
\begin{itemize}
  \item \texttt{client\_calls\_replay\_on\_reconnect\_buggy.tex}: a session's
    end leaves its calls pending and re-queues them for a live successor,
    the callback hold applied to requests.  C4 fails: a pending call outlives
    its session, and a request is served by a session it was not sent to.
  \item \texttt{client\_calls\_reply\_any\_session\_buggy.tex}: a reply is
    matched by request id alone.  A second live connection that is not the
    provider completes a call with a forged reply.  C2 fails.
  \item \texttt{client\_calls\_no\_deadline\_buggy.tex}: the guard of
    \texttt{Timeout} is false.  C5 fails.
  \item \texttt{client\_calls\_session\_end\_keeps\_pending\_buggy.tex}: the
    session ends and its calls stay pending until the deadline.  C4a fails,
    the same-step rule.
\end{itemize}

\end{document}
